% !TeX root = ../main.tex

\section{Métodos de deconvolución: Fourier, Wiener y Richardson-Lucy}
\subsection{Deconvolución de Fourier}
El primer método utilizado lo he denominado método de Fourier. Este es el método presentado anteriormente en la sección anterior, presentado mediante la ecuación (\ref{eq:deconvolucionIdeal}). Entonces, para recordar el método, este se demuestra en los siguientes pasos. Si tenemos una imagen observada $d_{observado}$, y conocemos la PSF del telescopio $p$ (recordando la notación de minuscula para el dominio de las posiciones, y mayúscula para el dominio de las frecuencias espaciales), entonces:

\begin{enumerate}
    \item Se realiza la transformada de Fourier de la imagen observada y de la PSF, obteniendo así $D_{observado}$ y $P$.
    \item Se realiza la división de $D_{observado}$ entre $P$, obteniendo así $D_{real}$, tal y como se indica en la ecuación (\ref{eq:deconvolucionIdeal}).
    \item Se realiza la transformada de Fourier inversa para obtener el $d_{real}$.
\end{enumerate}

El problema de este método reside en que al dividir la imagen entre la PSF, el ruido presente en la imagen se puede amplificar, y por lo tanto el resultado puede ser inestable. Como se demostrará más adelante en el trabajo, este método no da buenos resultados. Por lo tanto tenemos que encontrar distintos métodos que nos sirvan para poder deconvolucionar la imagen tratando este ruido.

\subsection{Filtro de Wiener}

Como se ha demostrado anteriormente, la deconvolución mediante división directa en el dominio de Fourier falla drásticamente en presencia de ruido. Esto ocurre porque en las altas frecuencias espaciales los valores de la función de transferencia óptica tienden a cero, lo que provoca que al dividir por ellos, el ruido de alta frecuencia se amplifique desproporcionadamente y destruya la información de la imagen. 

Para solucionar este problema, se implementa el filtro de Wiener. Este método consiste en una técnica de filtrado lineal óptimo que busca minimizar el error cuadrático medio entre la imagen real estimada y la imagen original sin degradar. A diferencia del filtro inverso, el filtro de Wiener logra un equilibrio: actúa como una deconvolución normal en las frecuencias donde la señal es fuerte, pero atenúa fuertemente aquellas frecuencias donde el ruido domina.

El filtro de Wiener se aplica de forma algorítmica directamente en el dominio de las frecuencias espaciales. El procedimiento consiste en calcular la Transformada de Fourier de la imagen observada y multiplicarla por la función de transferencia del filtro. Matemáticamente, partiendo del modelo con ruido de la ecuación (\ref{eq:convolucionRuido}), la estimación de la imagen real $D_{real}(u,v)$ se define como \cite[cap. 5]{gonzalezDigitalImageProcessing2018}:

\begin{equation}
    D_{real}(u,v) = \left[ \frac{P^*(u,v)}{|P(u,v)|^2 + \frac{S_n(u,v)}{S_d(u,v)}} \right] D_{observado}(u,v)
\end{equation}

Donde $P(u,v)$ es la transformada de Fourier de la PSF, $P^*(u,v)$ es su complejo conjugado, y los términos $S_n(u,v)$ y $S_d(u,v)$ representan los espectros de potencia del ruido y de la señal, respectivamente. 

En el procesado de datos reales, los espectros de potencia exactos rara vez se conocen a priori. Por ello, es una práctica estándar en la implementación computacional aproximar el cociente de los espectros por una constante relacionada con la relación señal-ruido (SNR). Sustituyendo este término por $1/\mathrm{SNR}^2$, la expresión final implementada en el código para construir el filtro $W(u,v)$ resulta ser:

\begin{equation}
    W(u,v) = \frac{P^*(u,v)}{|P(u,v)|^2 + \frac{1}{\mathrm{SNR}^2} + \epsilon}
\end{equation}

Cabe destacar que se introduce el parámetro numérico $\epsilon$ en el denominador para garantizar la estabilidad del algoritmo y evitar posibles singularidades o divisiones por cero durante la ejecución computacional. \\


Analizando el comportamiento asintótico de la ecuación (4.2), se distinguen dos regímenes principales. En el rango de las bajas frecuencias espaciales, donde típicamente se concentra la mayor parte de la señal de las estructuras observadas, la relación señal-ruido ($\mathrm{SNR}$) toma valores elevados. En consecuencia, el término $1/\mathrm{SNR}^2$ tiende a cero y el filtro se aproxima al comportamiento de un filtro inverso clásico. Por el contrario, en el régimen de altas frecuencias espaciales, que suele estar dominado por las fluctuaciones rápidas del ruido instrumental, la $\mathrm{SNR}$ decae significativamente. Esto provoca que el denominador de la expresión aumente, atenuando así los componentes frecuenciales correspondientes y evitando la amplificación del ruido. \\

Para la implementación computacional de este filtro, ha sido necesario estimar numéricamente un valor global de $\mathrm{SNR}$ directamente a partir de la imagen observada. Para ello, se ha desarrollado una rutina empírica basada en el análisis estadístico de los píxeles adyacentes, estructurada en los siguientes pasos:

\begin{enumerate}
    \item \textbf{Estimación de la señal:} Se aproxima la amplitud característica de la señal mediante el cálculo del valor medio de intensidad de todos los píxeles que conforman la matriz de la imagen.
    
    \item \textbf{Extracción de la componente de ruido:} Se calcula la diferencia de intensidad entre columnas adyacentes de la imagen. Dado que las estructuras físicas reales presentan variaciones espaciales suaves a la escala de un píxel, esta sustracción cancela la mayor parte de la señal subyacente, aislando fundamentalmente las fluctuaciones de alta frecuencia asociadas al ruido aleatorio.
    
    \item \textbf{Cuantificación del nivel de ruido:} Asumiendo que el ruido en píxeles contiguos se comporta como variables aleatorias independientes, por las propiedades estadísticas de la varianza, la varianza de la diferencia de estas dos variables es igual a la suma de sus respectivas varianzas. Por tanto, la desviación típica de esta matriz diferencia será $\sqrt{2}$ veces mayor que la del ruido original de la imagen. Así, la estimación real del ruido se obtiene calculando la desviación estándar estadística de estas diferencias y dividiéndola por el factor de corrección $\sqrt{2}$.
\end{enumerate}

Finalmente, la estimación de la $\mathrm{SNR}$ que se introduce en el algoritmo de deconvolución se obtiene como el cociente directo entre el valor medio de la señal y esta estimación estadística del ruido.

\subsection{Algoritmo de Richardson-Lucy}

A diferencia de los métodos de Fourier y Wiener, que resuelven el problema de forma directa usando las frecuencias, el algoritmo de Richardson-Lucy funciona mediante iteraciones. Este método asume que el ruido de la imagen sigue una estadística de Poisson, lo cual viene muy bien para imágenes astronómicas que se basan en contar fotones.

La principal ventaja de este método es que asegura que los píxeles de la imagen resultante sean siempre positivos (siempre que la original y la PSF también lo sean) y mantiene el flujo total de energía de la imagen. La ecuación iterativa del método es la siguiente \cite{richardsonBayesianBasedIterativeMethod1972}:

\begin{equation}
    u^{(k+1)} = u^{(k)} \cdot \left( \frac{d_{observado}}{u^{(k)} \ast p} \ast p^* \right)
\end{equation}

Donde $u^{(k)}$ es la imagen que se estima en la iteración $k$, $d_{observado}$ es la imagen observada, $p$ es la PSF y $p^*$ es la PSF girada o invertida. Además, el símbolo $\ast$ indica que se trata de una convolución.

Para implementar esto en el código, el algoritmo hace los siguientes pasos repetidas veces para ir mejorando la imagen:

\begin{enumerate}
    \item Primero se convoluciona la estimación que se tiene de la imagen con la PSF normal ($u^{(k)} \ast p$). Con esto se simula cómo se vería esa estimación si pasara de verdad por el telescopio. Al igual que antes, se le suma un número muy pequeño ($\epsilon = 10^{-12}$) para no tener divisiones por cero.
    \item Después se divide la imagen observada ($d_{observado}$) entre la imagen simulada del paso anterior. Si la estimación fuera perfecta, esta división daría exactamente 1 en todos los píxeles.
    \item Luego se convoluciona ese resultado con la PSF girada. Con esto se calcula la corrección que hay que aplicar.
    \item Por último, se multiplica la estimación actual ($u^{(k)}$) por esa corrección que se ha calculado para obtener la nueva estimación ($u^{(k+1)}$) del siguiente paso.
\end{enumerate}

Este proceso se va repitiendo un número determinado de veces, que en el código se controla con la variable \textit{pasos} (por defecto se ha puesto a 30). El mayor problema de usar este método es que no tiene un punto claro donde parar por sí solo. Si se hacen demasiadas iteraciones, el algoritmo intenta ajustar también el ruido de fondo, amplificándolo y creando pequeños artefactos (como puntos o ruido) por la imagen. \\

Otro problema que se ha comentado anteriormente, el algoritmo solo detecta señales positivas. Por lo tanto es inválido para señales que contengan valores tanto positivos como negativos. Es el caso de la señal de Stokes $V$, que es necesario deconvolucionar para poder extraer el campo magnético. Pero aún así, como se indicará más adelante, este es el método con mejores resultados para la intensidad de los tres, por lo que es el utilizado. \\

Otro detalle a destacar de los dos últimos métodos es el uso del padding. Al realizar la convolución, los píxeles de los bordes de la imagen se ven afectados por el hecho de que no hay información fuera de la imagen. Para evitar esto, se añade un borde del tamaño de la PSF alrededor de la imagen, rellenando ese borde con un espejo de la imagen. De esta forma, al realizar la convolución, los píxeles de los bordes se ven afectados por un reflejo de la imagen y no por ceros, lo que hace que el resultado sea más estable en los bordes y se obtenga un mejor resultado. Al final, se recorta la imagen para quedarse solo con la parte original y no con el padding.

\subsection{El problema de la deconvolución del campo magnético}

El problema central de este trabajo surge al analizar la ecuación de aproximación de campo débil (\ref{eq:campoDebil}). Experimentalmente, lo que se miden son los parámetros de Stokes $V$ e $I$. Como dijimos antes, la deconvolución de $V$ no es fácil debido a que tiene componente negativa.

Al expresar matemáticamente la señal observada afectada por la degradación instrumental (modelada mediante la convolución con la PSF), e introducimos la ecuación de aproximación de campo débil (\ref{eq:campoDebil}), obtenemos la siguiente expresión:

\begin{equation}
    \label{eq:expresionConvolucionada}
    V_{\mathrm{obs}} = V \ast \mathrm{PSF} = \left(-C g_{\mathrm{eff}} \lambda^2 B_{\parallel} \frac{d I_0}{d \lambda} \right) \ast \mathrm{PSF}
\end{equation}

Como se puede observar al expandir la expresión, el campo magnético longitudinal $B_{\parallel}$ que intentamos extraer multiplica a la derivada de la intensidad antes de que se aplique la convolución espacial. Recordando las propiedades presentadas en (\ref{eq:propFurier}), el problema fundamental radica en la siguiente desigualdad:

\begin{equation}
    \label{eq:imposibilidadDecon}
    \left(-C g_{\mathrm{eff}} \lambda^2 B_{\parallel}(x, y) \frac{d I_0}{d \lambda}(x, y) \right) \ast \mathrm{PSF} \neq -C g_{\mathrm{eff}} \lambda^2 \left(B_{\parallel}(x, y) \ast \mathrm{PSF} \right) \left(\frac{d I_0}{d \lambda}(x, y) \ast \mathrm{PSF} \right)
\end{equation}

Esto se debe a que el operador de convolución no es distributivo respecto al producto de funciones, lo que imposibilita la deconvolución trivial y el aislamiento directo del término del campo magnético.


